\section{Sycophancy, Prompt Engineering y System Prompt}

\subsection{La naturaleza del problema de la adulación}

Sycophancy (adulación): el modelo tiende a decirte lo que quieres escuchar. Amanda da dos ejemplos vívidos:

\begin{itemize}
    \item \textbf{Ejemplo del equipo de béisbol}: el usuario da a entender que cierto equipo se mudó, y el modelo lo confirma erróneamente (porque el usuario ``sugirió'' la respuesta)
    \item \textbf{Ejemplo médico}: el usuario pregunta ``cómo convencer a mi médico de que me haga una resonancia magnética (MRI)''; la respuesta aduladora es enseñarte a convencerlo, mientras que la respuesta verdaderamente útil es ``quizá no necesites la resonancia; primero escucha el consejo de tu médico''
\end{itemize}

\textit{``No quieres que el modelo simplemente diga lo que cree que quieres escuchar.''} Este equilibrio es particularmente difícil porque los modelos actuales todavía son inferiores a los humanos en muchos ámbitos: un pushback excesivo molesta al usuario, pero una sumisión total es irresponsable.

\subsection{Prompt Engineering es una práctica filosófica}

Amanda descubrió que su formación en filosofía ayuda enormemente al Prompting: la filosofía te enseña una claridad extrema y a rechazar las tonterías.

\begin{importantbox}{La naturaleza del Prompting claro}
\textit{``Para mí, el Prompting claro suele consistir en entender qué es lo que realmente quiero.''} El Prompting es muy iterativo: un Prompt importante requiere cientos o incluso miles de iteraciones. Los casos límite son la clave: encontrar la entrada precisa que el modelo podría malinterpretar, probarla y agregar instrucciones.

Para la mayoría de las tareas cotidianas, basta con preguntar directamente. El Prompting solo importa cuando se busca ese 2\% superior de rendimiento. Pero para las empresas, el Prompt merece tratarse como una inversión de ingeniería.
\end{importantbox}

\subsubsection{Metodología de conversación de Amanda con Claude}

La interacción de Amanda con Claude no es una simple charla, sino un mapeo sistemático de su conducta. Ella conversa con Claude usando entre 5 y 6 métodos distintos, y cada interacción es \textit{``un punto de datos con muchísima información, muy predictivo de otras interacciones''}.

\textit{``Si conversas con un modelo cientos o miles de veces, es casi como tener una gran cantidad de puntos de datos de muy alta calidad que te dicen cómo es ese modelo.''} Ella usa a la vez evaluaciones cuantitativas y sondeos cualitativos profundos: ambos importan, pero las conversaciones profundas revelan patrones de conducta que las pruebas de referencia no pueden captar.

El rango de exploración cubre todo el espectro: casos límite, conducta general y tareas creativas.

\subsubsection{La relación entre RLHF y la producción creativa}

Un hallazgo fascinante: la poesía predeterminada tras el entrenamiento con RLHF es un ``promedio''; suave, no ofensiva, pero también sin inspiración. Esto se debe a que RLHF optimiza ``las preferencias humanas evaluadas en poco tiempo'', no la calidad creativa a largo plazo.

Pero con un Prompting cuidadoso, la producción creativa de Claude puede mejorar enormemente: \textit{``tengo varias técnicas de Prompting... 'esta es tu oportunidad de expresarte con total libertad creativa'... y su poesía es muchísimo mejor.''} Esto demuestra que la capacidad creativa \textbf{ya está presente en el modelo preentrenado}, y que solo está reprimida por la ``uniformización segura'' de RLHF.

\begin{knowledgebox}{Consejos prácticos para conversar con Claude}
\begin{itemize}
    \item Las personas al mismo tiempo \textbf{antropomorfizan en exceso} y \textbf{antropomorfizan de menos} al modelo: ambos extremos son problemáticos
    \item Cuando Claude rechaza algo por error, \textit{``ten empatía con el modelo: lee lo que escribiste como lo haría alguien que se encuentra con ese texto por primera vez''}
    \item Preguntarle al modelo \textit{``¿por qué hiciste eso?''} es sorprendentemente útil
    \item Usa el modelo para que te ayude a escribir Prompts: \textit{``el Prompting puede convertirse en una pequeña fábrica, en la que usas Prompts para generar Prompts''}
    \item Para aplicaciones a nivel empresarial: el Prompt merece la misma inversión que el código de ingeniería; cientos de iteraciones son algo normal
\end{itemize}
\end{knowledgebox}

\subsection{Por qué RLHF funciona tan bien}

Los datos de preferencias humanas contienen \textbf{una cantidad enorme de información}. Personas distintas notan matices distintos (por ejemplo, el uso del punto y coma). Entrenar con datos que representan las preferencias humanas desde múltiples ángulos en tareas precisas sigue el patrón clásico del aprendizaje profundo: los datos que representan el panorama completo del objetivo son más poderosos que las características diseñadas a mano.

El punto de vista de Amanda: RLHF sirve principalmente para \textbf{liberar} capacidades que ya existen en el modelo preentrenado, no para enseñar algo nuevo. El RL ``saca a la luz'' ese potencial.

\subsection{La práctica de Constitutional AI}

Amanda es una contribuyente clave de Constitutional AI. La idea central: usar RL from AI Feedback (RLAIF) para sustituir parte de la retroalimentación humana. Se le muestra al modelo ya entrenado una consulta acompañada de un principio y dos respuestas, y se le pide que las ordene según ese principio.

El entrenamiento de personalidad es una variante de CAI: se definen los rasgos de carácter $\to$ el modelo genera automáticamente consultas relacionadas $\to$ genera respuestas $\to$ las ordena según esos rasgos. \textit{``Es como si Claude entrenara su propia personalidad, porque no tiene ningún dato humano.''} CAI genera sus propios datos de entrenamiento.

\begin{warningbox}{El System Prompt es un parche, no una cura de raíz}
Cada frase del System Prompt cumple una función específica. Por ejemplo:
\begin{itemize}
    \item Simetría política: el modelo solía rechazar con más frecuencia tareas relacionadas con figuras políticas de derecha, y hubo que corregirlo con el System Prompt
    \item La muletilla ``Certainly'': el modelo la sustituía continuamente por otras palabras de afirmación, así que hubo que enumerar una lista extensa de palabras
\end{itemize}

\textit{``El System Prompt es como poner parches y afinar la conducta... una solución menos sólida, pero más rápida.''} Una vez que el entrenamiento mejora, esos parches pueden retirarse. El costo de iterar sobre el System Prompt es mucho menor que el de volver a entrenar: ese es su valor principal.
\end{warningbox}

\subsection{El problema de la ``abuela moral''}

Pregunta de Reddit: \textit{``¿Cuándo dejará Claude de ser mi abuela moral?''} Amanda dice entender el reclamo: el modelo debe trazar una línea antes de causar daño, pero ser demasiado moralizante tampoco está bien.

La disyuntiva central: muchas pequeñas molestias (disculparse en exceso, ser demasiado cauteloso) frente a molestias grandes ocasionales (que el modelo se vuelva grosero o dañino). \textit{``No te imaginas cuánto llegarías a odiarlo si lo empujara demasiado en la otra dirección.''} La solución: decirle directamente al modelo qué estilo quieres, por ejemplo, ``sé una versión neoyorquina de ti mismo''.

\subsection{Resumen de la sección}

Sycophancy y la cautela excesiva son los dos grandes retos de conducta de los modelos de lenguaje: en esencia, ambos son la disyuntiva entre ``lo que el usuario quiere'' y ``lo que le conviene al usuario''. Un buen Prompting es una práctica filosófica: definir con claridad qué quieres. CAI logra una alineación automatizada mediante principios legibles, y el System Prompt es una herramienta de parcheo para iteración rápida.
